\section{Discussion}
\label{sec:discussion}

\subsection{Answer to Research Questions}~\label{sec:disc:rq}

We answer the research questions as follows:

\begin{itemize}
  \item \textbf{RQ1:} Our parallel congruence closure algorithm achieves a
  meaningful speedup on random benchmarks, a synthetic benchmark suite designed
  to stress test our algorithm, and a set of circuit equivalence benchmarks. We 
  scale well up to 32 cores (\as{give a number here}), but began to plateau beyond
  that, especially on smaller benchmarks.
  \item \textbf{RQ2:} Our algorithm tends to perform better on benchmarks with a
  large number of merges per round. Having more rounds of congruence closure
  does not necessarily lead to worse performance, but it does not lead to better
  performance either as we only parallelize within rounds.
  \item \textbf{RQ3:} Our algorithm provides a meaningful speedup over a
  sequential implementation on the circuit equivalence benchmarks. It is difficult
  to evaluate our algorithm on SMT or equality saturation benchmarks, as congruence
  closure is part of a larger system and thus difficult to isolate. We leave it to
  future work to integrate our algorithm into an SMT solver or equality saturation
  engine.
\end{itemize}

\subsection{Algorithm Modifications andFuture Work}~\label{sec:disc:limit-alg}

There are several extensions to the algorithm as presented that we would like
to address in future work. 

\textbf{Synchronous Round Structure.} Our synchronous round structure could be 
limited in terms of parallel scaling. In particular, the round structure means that we

inherently
limits parallel scaling. In particular, the instance size processed at each
round is small relative to the size of the problem for several benchmarks we
have examined (like the \texttt{egg} benchmarks). We will likely have to take a
more asynchronous approach to achieve better thread scaling; this would entail
generating new congruences from discovered equalities immediately instead of
waiting for the round to complete.

\textbf{Backtracking.} 
\as{Zak can you fill in how we could make our implementation backtrackable}

\textbf{Proof Production.} Our implementation does not produce \emph{proofs} or
\emph{witnesses} of the equalities it discovers. For instance if our algorithm
discovers that $x = y$, it cannot replay the chain of equalities and congruences
that led to this conclusion. This is important for a number of reasons: it is
necessary to justify correctness and it is used to produce conflict clauses in
SMT solvers.

Finally, we would like to support our implementation to incremental updates
(where equalities and new E-node may arrive dynamically). This is important for
both equality saturation and SMT solvers, where new equalities may be discovered
via for example quantifier instantiation.
